% !TEX root = userguide.tex

\def\headdate{12th November 2021}

\section{Macro elements and subdomain integration}
\label{sec:macroelements}

In this section some examples will be discussed
using the $P_1\text{-iso-}P_2$ and $Q_1\text{-iso-}Q_2$ shape functions and higher-order variants.
These shape functions are used together
with the element shapes \texttt{elshape=33,34,35,36} and \texttt{elshape=105,106,107}, as shown in Figure~\ref{fig:elements3} and Figure~\ref{fig:elements5}, respectively.
These 
shape functions have the same nodal layout as the regular higher-order elements. However, the interpolation for the elements is divided into linear, bilinear or trilinear subelements\footnote{Note, that the additional element diagonal (in red in
Figure~\ref{fig:elements3}) has been chosen to connect node 2 and 9. Since there are two other possible choices, the element will not be rotational invariant.}. On each subelement there is a linear ($P_1$) or bi-/trilinear ($Q_1$) interpolation.

As such, a mesh consisting of macro elements is only a mesh consisting of 
linear, bilinear and/or trilinear elements and therefore not particularly
interesting. However, the shape functions  $P_1\text{-iso-}P_2$ and
$Q_1\text{-iso-}Q_2$, for example, can replace the quadratic interpolation of the velocity
in the Taylor-Hood $P_2/P_1$ triangles and $Q_2/Q_1$ quads, respectively.
The resulting elements ($P_1\text{-iso-}P_2/P_1$ triangles, 
$Q_1\text{-iso-}Q_2/Q_1$ quadrilaterals) have been introduced in
\cite{Bercovier79} and fulfill
the LBB criterion. For some applications linear shape functions might be preferred
over higher-order ones, but this is beyond the scope of the present text.

In order to use the macro elements, the following must be taken into account:
\begin{itemize}
  \item Generate a mesh using the higher-order shape functions and change the
        element shape in the mesh manually in the main program to the
        macro element shape. For that, you can use \texttt{mesh\_change} (see Sec.~\ref{sec:meshgenextra}).
  \item Use the macro shape functions
        in the element subroutines instead of the usual higher-order ones.
  \item Change the volume and boundary integration to a subdomain integration\footnote{The
splitting into subdomains has been implemented in a general way and 
subdivision into $m^2$ triangles, quads and $m^3$ hexahedrons for
$m=1,2,\ldots$ is available. For tetrahedron only the split into eight subtetrahedrons is currently available.}
\end{itemize}

The following existing examples have been modified to show the use of
the macro elements:
\begin{center}
\begin{tabular}{lll}
  Main problem & shape function & Section of original problem \\\hline
  \texttt{poisson1a} & $Q_1\text{-iso-}Q_2$ & \ref{sec:poisson1} \\ 
  \texttt{poisson1b} & $P_1\text{-iso-}P_2$ & \ref{sec:poisson1} \\ 
  \texttt{stokes1a} & $Q_1\text{-iso-}Q_2$ & \ref{sec:driven_cavity2D} \\ 
  \texttt{stokes9a} & $Q_1\text{-iso-}Q_2$ & \ref{sec:samplingobjects} \\ 
  \texttt{stokes50} & $P_1\text{-iso-}P_2$ & \ref{sec:stokes3Ddrivencavity} \\ 
  \texttt{couette1a} & $Q_1\text{-iso-}Q_2$ & \ref{sec:couette} \\ 
  \texttt{couette1b} & $P_1\text{-iso-}P_2$ & \ref{sec:couette} \\ 
  \texttt{sphere9a} & $P_1\text{-iso-}P_2$ & \ref{sec:sphere2} \\
 \texttt{mixed\_stokes2} & $Q_1\text{-iso-}Q_p$ & \ref{sec:mixedcontraction} \\
\end{tabular}
\end{center}
More examples are available in the \texttt{drop}, \texttt{moving\_boundary} and \texttt{fountain\_flow} examples.


